\chapter{Martin Hairer --- Fields Medal 2014}\label{ch2014martinhairer:martin-hairer-fields-medal-2014-a-comprehensive-report-on-his-work-and-impact}

\begin{center}\rule{0.5\linewidth}{0.5pt}\end{center}

\textbf{Prepared:} July 2026 \textbf{Subject:} Sir Martin Hairer KBE FRS
(b. November 14, 1975, Geneva, Switzerland) \textbf{Position:} Professor
of Mathematics, EPFL and Imperial College London \textbf{Doctoral
Advisor:} Jean-Pierre Eckmann \textbf{Fields Medal Citation:}
\emph{``For outstanding contributions to the theory of stochastic
partial differential equations, and in particular for the creation of a
theory of regularity structures for such equations.''}

\begin{center}\rule{0.5\linewidth}{0.5pt}\end{center}

\section{Table of Contents}\label{ch2014martinhairer:table-of-contents}

\begin{enumerate}
\def\labelenumi{\arabic{enumi}.}
\tightlist
\item
 Introduction
\item
 Biographical Background
\item
 Stochastic Partial Differential Equations: Context and Challenges
\item
 The KPZ Equation: A Central Problem
\item
 Solving the KPZ Equation
\item
 Rough Path Theory: Foundations
\item
 Controlled Rough Paths
\item
 The Theory of Regularity Structures: Overview
\item
 Algebraic Framework of Regularity Structures
\item
 Analytic Framework: The Reconstruction Theorem
\item
 Renormalisation in Regularity Structures
\item
 The Dynamical \(\Phi^4_3\) Model
\item
 Other Singular SPDEs
\item
 Ergodicity and Hypoellipticity for Infinite-Dimensional Systems
\item
 The Navier-Stokes Equations and Stochastic Fluids
\item
 Numerical Methods and Software
\item
 Broader Impact and Influence
\item
 Awards and Recognition
\item
 Selected Major Publications
\item
 Ongoing Work and Future Directions
\item
 Conclusion
\item
 References and Further Reading
\end{enumerate}

\begin{center}\rule{0.5\linewidth}{0.5pt}\end{center}

\section{1. Introduction}\label{ch2014martinhairer:introduction}

Martin Hairer was awarded the Fields Medal in 2014 at the age of 38 for
his groundbreaking contributions to the theory of stochastic partial
differential equations (SPDEs), and in particular for creating the
theory of regularity structures. This theory provides a general
framework for solving a large class of singular stochastic PDEs ---
equations that arise naturally in statistical mechanics, quantum field
theory, and fluid dynamics, but whose solutions had been considered
mathematically ill-defined.

The central difficulty that Hairer addressed is that many important
SPDEs involve products of ``generalized functions'' (distributions) that
are too singular for their product to be defined in the classical sense.
Space-time white noise, for example, is so irregular that its values are
not functions but distributions. When such noise appears inside a
nonlinear term --- as in the KPZ equation or the dynamical \(\Phi^4_3\)
model --- the equation becomes ill-posed: the nonlinear term is a
product of distributions, which has no rigorous meaning.

Hairer's theory of regularity structures resolves this difficulty by
providing an algebraic and analytic framework that gives rigorous
meaning to these ill-posed equations. The theory can be seen as a
far-reaching extension of the theory of rough paths (pioneered by Terry
Lyons in the 1990s) from one-dimensional time to multi-dimensional
space-time, and from ordinary differential equations to partial
differential equations.

Hairer's contributions are characterized by notable technical
power combined with deep physical intuition. He has a rare ability to
identify the core mathematical structure underlying complex physical
problems and to develop the precise tools needed to analyze them. His
work has not only solved specific long-standing problems but has created
a general theory that will continue to be applied for decades to come.

This report provides a detailed account of Hairer's important
contributions, the mathematical context in which they arose, the
techniques he developed, and their lasting impact on mathematics and
physics.

\begin{center}\rule{0.5\linewidth}{0.5pt}\end{center}

\section{2. Biographical Background}\label{ch2014martinhairer:biographical-background}

\subsection{2.1 Early Life and
Education}\label{ch2014martinhairer:early-life-and-education}

Martin Hairer was born on November 14, 1975, in Geneva, Switzerland. His
father, Ernst Hairer, is a well-known numerical analyst and professor of
mathematics at the University of Geneva, specializing in the numerical
solution of ordinary differential equations. Growing up in an academic
environment, Martin was exposed to mathematics and computing from an
early age.

As a school student at Collège Claparède in Geneva, Hairer developed an
interest in audio programming. He created sound editing software that he
entered into a school science competition. This software eventually
became ``Amadeus Pro,'' a professional audio editing application for
macOS that Hairer continues to develop and maintain under his company
HairerSoft.

Hairer studied at the University of Geneva, earning a Bachelor of
Science in Mathematics (1998), a Master of Science in Physics (1998),
and a PhD in Physics (2001) under the supervision of Jean-Pierre
Eckmann. His doctoral thesis, ``Comportement Asymptotique d'Équations à
Dérivées Partielles Stochastiques'' (Asymptotic Behavior of Stochastic
Partial Differential Equations), foreshadowed the direction of his
future work.

\subsection{2.2 Academic Career}\label{ch2014martinhairer:academic-career}

After completing his PhD, Hairer held positions at the University of
Warwick, the Courant Institute of Mathematical Sciences at NYU, and
Imperial College London. He was appointed Regius Professor of
Mathematics at the University of Warwick in 2014, a position he held
until 2017. He is currently Professor of Mathematics at both EPFL (École
Polytechnique Fédérale de Lausanne) and Imperial College London.

Hairer was knighted as a Knight Commander of the Order of the British
Empire (KBE) in 2025 for his services to mathematics. He is a Fellow of
the Royal Society and a member of numerous international academies.

\subsection{2.3 Personal Life}\label{ch2014martinhairer:personal-life}

Hairer is married to Xue-Mei Li, a distinguished mathematician
specializing in stochastic analysis and differential geometry who
currently holds a position at Imperial College London.

\begin{center}\rule{0.5\linewidth}{0.5pt}\end{center}

\section{3. Stochastic Partial Differential Equations: Context and
Challenges}\label{ch2014martinhairer:stochastic-partial-differential-equations-context-and-challenges}

\subsection{3.1 What Are SPDEs?}\label{ch2014martinhairer:what-are-spdes}

A stochastic partial differential equation (SPDE) is a partial
differential equation that contains a random noise term. SPDEs arise
naturally when modeling systems that are subject to random fluctuations
--- for example, the random motion of particles in a fluid, the random
growth of an interface, or the quantum fluctuations of a field.

A typical SPDE takes the form:

\[
\partial_t u = \Delta u + F(u) + \xi
\]

where \(\xi\) denotes ``space-time white noise'' --- a random field that
is delta-correlated in both space and time.

\subsection{3.2 The Problem of
Ill-Posedness}\label{ch2014martinhairer:the-problem-of-ill-posedness}

The fundamental difficulty with singular SPDEs is that space-time white
noise \(\xi\) is not a function but a distribution: it is too irregular
to be evaluated at individual points. In one spatial dimension, it is
just regular enough that the linear equation
\(\partial_t u = \Delta u + \xi\) has a well-defined solution (the
stochastic heat equation). However, in two or three spatial dimensions,
the noise is so irregular that nonlinear terms like \(F(u) = u^3\) or
\((\partial_x h)^2\) involve products of distributions, which are not
classically defined.

For decades, mathematicians struggled to give rigorous meaning to these
equations. The problem was not merely technical: it went to the heart of
what it means to be a solution to a differential equation.

\subsection{3.3 The Need for
Renormalisation}\label{ch2014martinhairer:the-need-for-renormalisation}

Physicists had long known that these equations require
``renormalisation'' --- the subtraction of infinite constants to obtain
finite, physically meaningful quantities. This is analogous to the
renormalisation program in quantum field theory. But a rigorous
mathematical framework for renormalisation in SPDEs was lacking before
Hairer's work.

\begin{center}\rule{0.5\linewidth}{0.5pt}\end{center}

\section{4. The KPZ Equation: A Central
Problem}\label{ch2014martinhairer:the-kpz-equation-a-central-problem}

\subsection{4.1 Physical Origin}\label{ch2014martinhairer:physical-origin}

The KPZ equation, named after Kardar, Parisi, and Zhang (1986),
describes the stochastic growth of an interface:

\[
\partial_t h = \partial_x^2 h + (\partial_x h)^2 + \xi
\]

Here \(h(t,x)\) is the height of the interface at position \(x\) and
time \(t\), and \(\xi\) is space-time white noise. The nonlinear term
\((\partial_x h)^2\) captures the tendency of the interface to grow
perpendicular to itself.

\subsection{4.2 Universality}\label{ch2014martinhairer:universality}

The KPZ equation is central to statistical mechanics because it belongs
to a universality class that describes a wide range of growth phenomena,
including: - Deposition of particles on a surface - Growth of bacterial
colonies - Propagation of flame fronts - Fluctuations in liquid crystals

The KPZ universality class is characterized by specific scaling
exponents (the roughness exponent \(\chi = 1/2\) and the dynamic
exponent \(z = 3/2\)) that are independent of the microscopic details of
the system.

\subsection{4.3 Mathematical Challenges}\label{ch2014martinhairer:mathematical-challenges}

Despite its physical importance, the KPZ equation was mathematically
ill-posed. The nonlinear term \((\partial_x h)^2\) involves the square
of a derivative of the solution, which is not a function but a
distribution when the noise is white. The equation requires
renormalisation: formally, one must subtract an infinite constant from
the nonlinear term to obtain a finite limit.

Before Hairer's work, the only way to define solutions to the KPZ
equation was through the Hopf-Cole transformation: setting \(Z = e^h\)
transforms the KPZ equation into the stochastic heat equation
\(\partial_t Z = \partial_x^2 Z + Z\xi\), which is linear and can be
solved. However, this trick works only for the specific case of the KPZ
equation with additive noise; it fails for generalizations such as
coupled KPZ systems or KPZ with multiplicative noise.

\begin{center}\rule{0.5\linewidth}{0.5pt}\end{center}

\section{5. Solving the KPZ Equation}\label{ch2014martinhairer:solving-the-kpz-equation}

\subsection{5.1 Hairer's 2013 Solution}\label{ch2014martinhairer:hairers-2013-solution}

In 2013, Hairer published a landmark paper ``Solving the KPZ equation''
in the Annals of Mathematics. Using ideas from rough path theory, he
gave the first direct mathematical construction of solutions to the KPZ
equation --- without relying on the Hopf-Cole transformation.

The key innovation was to treat the nonlinear term \((\partial_x h)^2\)
not as a product of distributions but as a ``rough path'' lift of the
solution. Hairer showed that, just as in rough path theory, one needs to
keep track of additional information --- in this case, the iterated
integrals of the noise --- to make sense of the equation.

\subsection{5.2 The Renormalisation
Procedure}\label{ch2014martinhairer:the-renormalisation-procedure}

Hairer's construction involved a renormalisation procedure:
approximating the noise by smooth functions (regularisation), solving
the regularised equation, subtracting a diverging constant from the
nonlinear term, and then taking the limit as the regularisation is
removed. He proved that this limit exists and is independent of the
approximation scheme, giving a well-defined notion of solution.

\subsection{5.3 The Cole-Hopf
Connection}\label{ch2014martinhairer:the-cole-hopf-connection}

Hairer also proved that his constructed solutions coincide with the
Hopf-Cole solutions when the latter exist. This confirmed that the
Hopf-Cole transformation was not just a trick but captured the correct
physics.

\subsection{5.4 Beyond Hopf-Cole}\label{ch2014martinhairer:beyond-hopf-cole}

Crucially, Hairer's method extended to equations where the Hopf-Cole
transformation is unavailable, such as coupled KPZ systems and KPZ with
multiplicative noise. This generalization was essential for applications
to physical systems with multiple interacting interfaces.

\begin{center}\rule{0.5\linewidth}{0.5pt}\end{center}

\section{6. Rough Path Theory:
Foundations}\label{ch2014martinhairer:rough-path-theory-foundations}

\subsection{6.1 Lyons's Theory}\label{ch2014martinhairer:lyonss-theory}

Rough path theory, developed by Terry Lyons in the 1990s, provides a
pathwise theory of differential equations driven by irregular signals.
The central idea is that to make sense of the integral \(\int Y dX\)
when \(X\) is rough (e.g., a Brownian path), one needs to augment \(X\)
with its iterated integrals --- the ``rough path'' lift:

\[
\mathbb{X}_{s,t} = \left(1, \int_s^t dX_{r_1}, \int_s^t \int_s^{r_2} dX_{r_1} \otimes dX_{r_2}, \ldots\right)
\]

\subsection{6.2 The Signature}\label{ch2014martinhairer:the-signature}

The collection of all iterated integrals of a path is called its
``signature.'' The signature encodes the path in a way that is
sufficient for solving differential equations: if one knows the
signature of \(X\) to a sufficiently high order, one can define
\(\int Y dX\) for any \(Y\) that is itself controlled by \(X\).

\subsection{6.3 Limitations}\label{ch2014martinhairer:limitations}

Rough path theory was restricted to one-dimensional time. It could
handle stochastic ordinary differential equations (SODEs) but not
stochastic partial differential equations (SPDEs), where the noise
varies in both space and time.

\begin{center}\rule{0.5\linewidth}{0.5pt}\end{center}

\section{7. Controlled Rough Paths}\label{ch2014martinhairer:controlled-rough-paths}

\subsection{7.1 Gubinelli's Controlled
Paths}\label{ch2014martinhairer:gubinellis-controlled-paths}

Hairer, together with Massimiliano Gubinelli, developed the theory of
``controlled rough paths,'' which provides a more flexible framework for
rough path analysis. Instead of requiring the integrand \(Y\) to be a
function of the path \(X\), controlled rough path theory allows \(Y\) to
be any process that is ``controlled'' by \(X\) in the sense that its
increments can be expressed in terms of \(X\)'s increments plus a
higher-order remainder.

\subsection{7.2 The Space of Controlled
Paths}\label{ch2014martinhairer:the-space-of-controlled-paths}

A path \(Y\) is said to be controlled by \(X\) if there exists a process
\(Y'\) such that:

\[
Y_t - Y_s = Y'_s (X_t - X_s) + O(|t-s|^{2\alpha})
\]

for some \(\alpha > 1/3\). This condition ensures that the integral
\(\int Y dX\) can be defined using the rough path lift of \(X\).

\subsection{7.3 Application to SPDEs}\label{ch2014martinhairer:application-to-spdes}

The theory of controlled rough paths was crucial for Hairer's solution
of the KPZ equation. By treating the solution \(h\) as a function
controlled by the noise, Hairer was able to define the problematic
nonlinear term \((\partial_x h)^2\) without requiring additional
regularity.

\begin{center}\rule{0.5\linewidth}{0.5pt}\end{center}

\section{8. The Theory of Regularity Structures:
Overview}\label{ch2014martinhairer:the-theory-of-regularity-structures-overview}

\subsection{8.1 Motivation}\label{ch2014martinhairer:motivation}

While rough path theory could handle the one-dimensional KPZ equation
(one spatial dimension + one time dimension), higher-dimensional SPDEs
required a substantially new approach. The theory of regularity
structures, introduced by Hairer in his 2014 paper ``A theory of
regularity structures'' (published in Inventiones Mathematicae),
provides this generalization.

\subsection{8.2 The Core Idea}\label{ch2014martinhairer:the-core-idea}

Regularity structures generalize the concept of a Taylor expansion. In
classical analysis, a smooth function can be approximated locally by a
polynomial (its Taylor polynomial). The theory of regularity structures
replaces polynomials with more general ``models'' that are purpose-built
for the equation at hand.

A regularity structure consists of three components: 1. \textbf{An index
set} \(A \subset \mathbb{R}\) (the set of ``regularities'') 2. \textbf{A
graded vector space} \(T = \oplus_{\alpha \in A} T_\alpha\) (the space
of ``modelled distributions'') 3. \textbf{A structure group} \(G\) (the
group of transformations that relate different local expansions)

\subsection{8.3 The Reconstruction
Theorem}\label{ch2014martinhairer:the-reconstruction-theorem}

The centerpiece of the theory is the Reconstruction Theorem, which
states that a ``model'' (a consistent family of local approximations)
can be ``reconstructed'' into a global distribution. This theorem is the
analogue of the fact that a smooth function is determined by its Taylor
expansions at every point.

\subsection{8.4 Application to SPDEs}\label{ch2014martinhairer:application-to-spdes-1}

To solve a singular SPDE using regularity structures, one: 1. Identifies
the relevant regularity structure (the set of ``symbols'' needed to
describe the solution) 2. Constructs a model for this structure
(encoding the noise and its iterated integrals) 3. Solves the equation
as a fixed point in the space of modelled distributions 4. Proves that
the solution converges as the regularization of the noise is removed

\begin{center}\rule{0.5\linewidth}{0.5pt}\end{center}

\section{9. Algebraic Framework of Regularity
Structures}\label{ch2014martinhairer:algebraic-framework-of-regularity-structures}

\subsection{9.1 Trees and Symbols}\label{ch2014martinhairer:trees-and-symbols}

A key feature of regularity structures is the use of combinatorial
objects --- rooted trees decorated with labels --- to represent the
terms that appear in the expansion of the solution. Each tree
corresponds to a ``symbol'' that captures a particular combination of
noise terms and their derivatives.

For example, in the \(\Phi^4_3\) equation, the symbols include: -
\(\Xi\): the noise - \(\mathcal{I}(\Xi)\): the heat kernel applied to
the noise - \(\mathcal{I}(\Xi^3)\): the heat kernel applied to the cube
of the noise - More complex trees representing the renormalisation
counterterms

\subsection{9.2 The Structure Group}\label{ch2014martinhairer:the-structure-group}

The structure group \(G\) acts on the space of symbols by
``relabelling'' the expansions. This group encodes the symmetries of the
equation and plays a crucial role in the renormalisation procedure.

\subsection{9.3 The Model Space}\label{ch2014martinhairer:the-model-space}

The model space \(T\) is the graded vector space spanned by the symbols.
The grading is given by the ``homogeneity'' of each symbol, which
determines its regularity. The solution to the SPDE is represented as a
modelled distribution --- a function that takes values in \(T\) and
satisfies certain analytic constraints.

\begin{center}\rule{0.5\linewidth}{0.5pt}\end{center}

\section{10. Analytic Framework: The Reconstruction
Theorem}\label{ch2014martinhairer:analytic-framework-the-reconstruction-theorem}

\subsection{10.1 Statement}\label{ch2014martinhairer:statement}

Let \(\mathcal{T} = (A, T, G)\) be a regularity structure. A
\textbf{model} for \(\mathcal{T}\) consists of: - A family of linear
maps \(\Pi_x: T \to \mathcal{S}'(\mathbb{R}^d)\) (the ``realizations''
of symbols at each point \(x\)) - A family of linear maps
\(\Gamma_{xy}: T \to T\) (the ``connection maps'' relating different
points)

satisfying certain compatibility conditions.

The \textbf{Reconstruction Theorem} states that for any modelled
distribution \(f\) (a function taking values in \(T\) that satisfies a
Hölder-type condition), there exists a unique distribution
\(\mathcal{R}f\) such that:

\[
\mathcal{R}f(x) \approx \Pi_x f(x)
\]

in a precise sense. This theorem is the analytic foundation of the
theory.

\subsection{10.2 Schauder Estimates}\label{ch2014martinhairer:schauder-estimates}

Hairer proved Schauder-type estimates for regularity structures: if a
modelled distribution has regularity \(\gamma\), then its reconstruction
has regularity \(\alpha = \gamma + \beta\) (where \(\beta\) is the
regularity of the model). These estimates are essential for solving
SPDEs by fixed-point arguments.

\subsection{10.3 Fixed-Point Solutions}\label{ch2014martinhairer:fixed-point-solutions}

To solve an SPDE such as \(\partial_t u = \Delta u + F(u) + \xi\) using
regularity structures: 1. Write the equation in mild form:
\(u = K * (F(u) + \xi)\), where \(K\) is the heat kernel 2. Represent
\(u\) as a modelled distribution in an appropriate regularity structure
3. Show that the map \(u \mapsto K * (F(u) + \xi)\) is a contraction on
a suitable space of modelled distributions 4. Apply the Banach
fixed-point theorem

\begin{center}\rule{0.5\linewidth}{0.5pt}\end{center}

\section{11. Renormalisation in Regularity
Structures}\label{ch2014martinhairer:renormalisation-in-regularity-structures}

\subsection{11.1 The Need for
Renormalisation}\label{ch2014martinhairer:the-need-for-renormalisation-1}

In singular SPDEs, certain ``Feynman diagrams'' (represented by trees in
the regularity structure) become divergent because they involve products
of distributions that are not well-defined. These divergences manifest
as infinite constants that must be subtracted to obtain finite
solutions.

\subsection{11.2 The Renormalisation
Group}\label{ch2014martinhairer:the-renormalisation-group}

Hairer introduced a renormalisation group within the framework of
regularity structures. The group consists of transformations of the
model that correspond to adding counterterms to the equation. By
applying appropriate transformations, one can cancel the divergences and
obtain a well-defined limit.

\subsection{11.3 The Renormalised Model}\label{ch2014martinhairer:the-renormalised-model}

The key result is that there exists a renormalised model --- a modified
version of the original model --- for which the fixed-point equation has
a unique solution. This renormalised model is obtained by composing the
original model with an element of the renormalisation group, which
corresponds to subtracting diverging constants from the nonlinear terms.

\subsection{11.4 Universality}\label{ch2014martinhairer:universality-1}

Hairer proved that the renormalised solutions are independent of the
regularization scheme: any reasonable approximation of the noise leads
to the same limit after renormalisation. This universality property is
essential for the physical interpretation of the solutions.

\begin{center}\rule{0.5\linewidth}{0.5pt}\end{center}

\section{\texorpdfstring{12. The Dynamical \(\Phi^4_3\)
Model}{12. The Dynamical \textbackslash Phi\^{}4\_3 Model}}\label{ch2014martinhairer:the-dynamical-phi4_3-model}

\subsection{12.1 The Equation}\label{ch2014martinhairer:the-equation}

The dynamical \(\Phi^4_3\) model is the stochastic quantization equation
for the Euclidean \(\Phi^4\) quantum field theory in three spatial
dimensions:

\[
\partial_t \Phi = \Delta \Phi + C \Phi - \Phi^3 + \xi
\]

where \(\xi\) is space-time white noise, \(C\) is a renormalisation
constant that diverges as the regularization is removed, and the spatial
domain is a three-dimensional torus.

This equation describes the stochastic evolution of a scalar field
\(\Phi(t,x)\) with a quartic self-interaction. It arises in quantum
field theory as the Langevin equation for the \(\Phi^4\) Euclidean
quantum field.

\subsection{12.2 Previous Work}\label{ch2014martinhairer:previous-work}

Before Hairer's work, the \(\Phi^4_3\) model had been studied using
constructive quantum field theory methods, which established the
existence of the stationary measure (the \(\Phi^4_3\) Euclidean quantum
field) but did not provide a dynamical theory. The stochastic
quantization approach --- which would give the dynamics --- was
considered mathematically intractable because the equation is ill-posed.

\subsection{12.3 Hairer's Construction}\label{ch2014martinhairer:hairers-construction}

Using regularity structures, Hairer constructed solutions to the
\(\Phi^4_3\) equation for spatial dimensions \(d < 4\). The construction
involves: 1. Identifying the relevant trees (symbols) up to a certain
degree 2. Computing the divergences associated with each tree 3.
Defining the renormalised model by subtracting the divergences 4.
Solving the fixed-point equation in the space of modelled distributions

The result is a local-in-time solution that is almost surely continuous
in time and takes values in a space of distributions.

\subsection{12.4 The Critical Dimension}\label{ch2014martinhairer:the-critical-dimension}

The dimension \(d = 4\) is critical for the \(\Phi^4\) model: above this
dimension, the theory is trivial (non-interacting); below this
dimension, it is super-renormalisable. Hairer's theory covers dimensions
2 and 3, which are the physically most interesting cases.

\begin{center}\rule{0.5\linewidth}{0.5pt}\end{center}

\section{13. Other Singular SPDEs}\label{ch2014martinhairer:other-singular-spdes}

\subsection{13.1 The Parabolic Anderson
Model}\label{ch2014martinhairer:the-parabolic-anderson-model}

The parabolic Anderson model:

\[
\partial_t u = \Delta u + u \xi
\]

describes the evolution of a particle in a random potential. In two and
three dimensions, this equation requires renormalisation because the
product \(u\xi\) involves the multiplication of distributions.

Using regularity structures, Hairer and collaborators constructed
solutions to the parabolic Anderson model in dimensions 2 and 3,
providing the first rigorous treatment of this equation.

\subsection{13.2 Stochastic Navier-Stokes
Equations}\label{ch2014martinhairer:stochastic-navier-stokes-equations}

In joint work with Jonathan Mattingly, Hairer studied the stochastic
Navier-Stokes equations with highly singular forcing. While the
deterministic Navier-Stokes equations are well understood in two
dimensions, adding space-time white noise creates significant
difficulties. Hairer's work on hypoellipticity (discussed below)
provides tools for analyzing the ergodic properties of these equations.

\subsection{13.3 Generalized KPZ
Equations}\label{ch2014martinhairer:generalized-kpz-equations}

Hairer extended his KPZ results to a broader class of equations of the
form:

\[
\partial_t h = \partial_x^2 h + g_1(h)(\partial_x h)^2 + g_2(h)\partial_x h + g_3(h) + g_4(h)\xi
\]

These generalized KPZ equations arise in the study of fluctuating
interfaces with coordinate-dependent growth rates. They cannot be solved
by the Hopf-Cole transformation because the transformation does not
eliminate the nonlinearity.

\subsection{13.4 Coupled Systems}\label{ch2014martinhairer:coupled-systems}

Hairer's theory also handles systems of coupled SPDEs, such as multiple
KPZ equations coupled through their nonlinearities. Such systems appear
in the study of chains of nonlinearly interacting oscillators and in
models of universality classes beyond KPZ.

\begin{center}\rule{0.5\linewidth}{0.5pt}\end{center}

\section{14. Ergodicity and Hypoellipticity for Infinite-Dimensional
Systems}\label{ch2014martinhairer:ergodicity-and-hypoellipticity-for-infinite-dimensional-systems}

\subsection{14.1 The Problem of
Ergodicity}\label{ch2014martinhairer:the-problem-of-ergodicity}

A fundamental question in stochastic dynamics is whether the system has
a unique stationary measure and converges to it exponentially fast
(ergodicity). For SPDEs, this question is nontrivial because the state
space is infinite-dimensional.

\subsection{14.2 Hairer's Hypoellipticity
Theory}\label{ch2014martinhairer:hairers-hypoellipticity-theory}

In a series of papers with Jonathan Mattingly, Hairer developed a
general theory of hypoellipticity for infinite-dimensional stochastic
systems. Hypoellipticity refers to the smoothing property of stochastic
differential equations: even if the noise only acts directly on some
components of the system, the dynamics may still be smooth in all
directions.

\subsection{14.3 Key Results}\label{ch2014martinhairer:key-results}

Hairer and Mattingly proved: - Unique ergodicity for the stochastic
Navier-Stokes equations with degenerate noise - Exponential convergence
to the stationary measure - A criterion for hypoellipticity in infinite
dimensions based on the ``bracket condition'' (the stochastic analogue
of Hörmander's condition)

These results are among the most general available for ergodicity of
SPDEs.

\begin{center}\rule{0.5\linewidth}{0.5pt}\end{center}

\section{15. The Navier-Stokes Equations and Stochastic
Fluids}\label{ch2014martinhairer:the-navier-stokes-equations-and-stochastic-fluids}

\subsection{15.1 The Stochastic Navier-Stokes
Equations}\label{ch2014martinhairer:the-stochastic-navier-stokes-equations}

The Navier-Stokes equations describe the motion of viscous fluids.
Adding stochastic forcing models the effects of random external forces
or thermal fluctuations:

\[
\partial_t v = \Delta v - (v \cdot \nabla)v - \nabla p + \xi, \quad \nabla \cdot v = 0
\]

\subsection{15.2 Hairer's Contributions}\label{ch2014martinhairer:hairers-contributions}

Hairer's work on the stochastic Navier-Stokes equations includes: -
Proof of unique ergodicity for the 2D stochastic Navier-Stokes equations
with degenerate noise - Development of general criteria for ergodicity
that apply beyond the Navier-Stokes equations - Analysis of the
multiscale behavior of solutions

\subsection{15.3 The ``Hypoelliptic
Navier-Stokes''}\label{ch2014martinhairer:the-hypoelliptic-navier-stokes}

The key difficulty is that the noise typically acts only on a finite
number of Fourier modes, yet the dynamics must propagate this randomness
to all modes. Hairer and Mattingly showed that the nonlinear term
\((v \cdot \nabla)v\) provides the necessary ``hypoelliptic'' structure
to spread the noise, resulting in a unique stationary measure.

\begin{center}\rule{0.5\linewidth}{0.5pt}\end{center}

\section{16. Broader Impact and
Influence}\label{ch2014martinhairer:broader-impact-and-influence}

\subsection{16.1 The Creation of a New
Field}\label{ch2014martinhairer:the-creation-of-a-new-field}

Hairer's theory of regularity structures has created a new subfield of
stochastic analysis. It has been taken up by researchers around the
world and applied to problems ranging from quantum field theory to
finance.

\subsection{16.2 Regularity Structures and Paracontrolled
Distributions}\label{ch2014martinhairer:regularity-structures-and-paracontrolled-distributions}

Around the same time as Hairer's work, Gubinelli, Imkeller, and
Perkowski developed the theory of ``paracontrolled distributions,''
which provides an alternative approach to singular SPDEs using
paraproducts from harmonic analysis. The two theories are complementary:
paracontrolled distributions are simpler in some respects but less
general than regularity structures.

\subsection{16.3 Applications to Quantum Field
Theory}\label{ch2014martinhairer:applications-to-quantum-field-theory}

Regularity structures have found important applications in constructive
quantum field theory. The stochastic quantization approach to Euclidean
quantum field theory --- pioneered by Parisi and Wu in 1981 --- uses
SPDEs to construct quantum fields. Hairer's work provides the rigorous
foundation for this approach in dimensions 2 and 3.

\subsection{16.4 Rough Volatility in
Finance}\label{ch2014martinhairer:rough-volatility-in-finance}

In quantitative finance, ``rough volatility'' models --- which use
fractional Brownian motion with small Hurst parameter to model
volatility --- have become increasingly important. The theory of rough
paths and regularity structures provides the mathematical tools needed
to analyze these models.

\subsection{16.5 Influence on Young
Researchers}\label{ch2014martinhairer:influence-on-young-researchers}

Hairer has supervised numerous PhD students and postdocs who have gone
on to become leaders in stochastic analysis. His monograph ``A Course on
Rough Paths'' (with Peter Friz) has become the standard introduction to
the subject.

\subsection{16.6 The Amadeus Software}\label{ch2014martinhairer:the-amadeus-software}

Beyond his mathematical work, Hairer continues to develop Amadeus Pro, a
professional audio editing application for macOS. The software, which
originated from his school science project, is widely used for audio
editing, analysis, and processing. This combination of pure mathematics
with practical software development is unusual and reflects Hairer's
diverse talents.

\begin{center}\rule{0.5\linewidth}{0.5pt}\end{center}

\section{17. Awards and Recognition}\label{ch2014martinhairer:awards-and-recognition}

\textbf{2025:} - Sylvester Medal (Royal Society) - Knight Commander of
the Order of the British Empire (KBE)

\textbf{2022:} - King Faisal Prize in Mathematics - ESI Medal (Erwin
Schrödinger International Institute)

\textbf{2021:} - Breakthrough Prize in Mathematics

\textbf{2014:} - Fields Medal (International Mathematical Union) -
Fröhlich Prize (London Mathematical Society)

\textbf{2013:} - Fermat Prize (Institut de Mathématiques de Toulouse)

\textbf{2009:} - Wolfson Research Merit Award (Royal Society)

\textbf{2008:} - Whitehead Prize (London Mathematical Society) - Philip
Leverhulme Prize

\textbf{2006:} - Royal Society University Research Fellowship

\begin{center}\rule{0.5\linewidth}{0.5pt}\end{center}

\section{18. Selected Major
Publications}\label{ch2014martinhairer:selected-major-publications}

\begin{enumerate}
\def\labelenumi{\arabic{enumi}.}
\item
 \textbf{``A theory of regularity structures.''} \emph{Inventiones
 Mathematicae} 198(2), 269--504, 2014. --- \emph{The foundational paper
 introducing regularity structures, providing a general framework for
 solving singular SPDEs.}
\item
 \textbf{``Solving the KPZ equation.''} \emph{Annals of Mathematics}
 178(2), 559--664, 2013. --- \emph{First direct construction of
 solutions to the KPZ equation using rough path methods.}
\item
 \textbf{``Ergodicity for the stochastic Navier-Stokes equations with
 degenerate noise''} (with J. Mattingly). \emph{Annals of Mathematics}
 164(3), 993--1032, 2006. --- \emph{Proof of unique ergodicity for 2D
 stochastic Navier-Stokes with degenerate forcing.}
\item
 \textbf{``Hypoellipticity in infinite dimensions''} (with J.
 Mattingly). \emph{Communications on Pure and Applied Mathematics}
 64(6), 765--828, 2011. --- \emph{General theory of hypoellipticity for
 infinite-dimensional stochastic systems.}
\item
 \textbf{``The strong Feller property for the stochastic Navier-Stokes
 equation''} (with J. Mattingly). \emph{Journal of Functional Analysis}
 239(2), 363--407, 2006. --- \emph{Smoothing properties of the
 stochastic Navier-Stokes semigroup.}
\item
 \textbf{``Rough stochastic PDEs.''} \emph{Communications on Pure and
 Applied Mathematics} 64(11), 1547--1585, 2011. --- \emph{Extension of
 rough path theory to SPDEs with one-dimensional spatial domain.}
\item
 \textbf{``Controlled rough paths''} (with M. Gubinelli). \emph{Journal
 of Functional Analysis} 266(3), 1205--1235, 2014. ---
 \emph{Development of the theory of controlled rough paths.}
\item
 \textbf{``A Course on Rough Paths''} (with P. Friz). Springer, 2014
 (2nd ed.~2020). --- \emph{Standard textbook on rough path theory and
 regularity structures.}
\item
 \textbf{``Introduction to Regularity Structures.''} \emph{Brazilian
 Journal of Probability and Statistics} 29(2), 175--210, 2015. ---
 \emph{Overview of regularity structures for a general audience.}
\item
 \textbf{``Singular stochastic PDEs.''} \emph{Proceedings of the
 International Congress of Mathematicians} 2014. --- \emph{ICM survey
 of Hairer's work on singular SPDEs.}
\end{enumerate}

\begin{center}\rule{0.5\linewidth}{0.5pt}\end{center}

\section{19. Ongoing Work and Future
Directions}\label{ch2014martinhairer:ongoing-work-and-future-directions}

\subsection{19.1 Extensions of Regularity
Structures}\label{ch2014martinhairer:extensions-of-regularity-structures}

Hairer and collaborators continue to extend the theory of regularity
structures to new classes of equations: - Multi-scale regularity
structures for systems with multiple characteristic scales - Regularity
structures for non-parabolic equations (e.g., hyperbolic SPDEs) -
Regularity structures for equations with non-local operators

\subsection{19.2 Quantum Field Theory}\label{ch2014martinhairer:quantum-field-theory}

The application of regularity structures to quantum field theory is an
active area of research. Hairer is a principal investigator in the
Simons Collaboration on Probabilistic Paths to Quantum Field Theory,
which aims to use stochastic methods to construct non-perturbative
quantum field theories.

\subsection{19.3 Numerical Methods}\label{ch2014martinhairer:numerical-methods}

The theory of regularity structures suggests new numerical methods for
simulating singular SPDEs. By incorporating the algebraic structure of
the theory, one can design numerical schemes that correctly capture the
renormalisation needed for stable simulations.

\subsection{19.4 Machine Learning and Rough
Paths}\label{ch2014martinhairer:machine-learning-and-rough-paths}

Recent work has explored connections between rough path theory and
machine learning. The signature of a path provides a natural set of
features for learning tasks involving sequential data, and regularity
structures may provide a framework for learning solutions to PDEs.

\begin{center}\rule{0.5\linewidth}{0.5pt}\end{center}

\section{20. Hairer's Mathematical
Style}\label{ch2014martinhairer:hairers-mathematical-style}

\subsection{20.1 Physical Intuition}\label{ch2014martinhairer:physical-intuition}

Hairer's work is distinguished by its deep physical intuition. He
approaches SPDEs not as abstract mathematical objects but as
descriptions of physical phenomena, and his constructions are guided by
physical principles such as universality and renormalisation.

\subsection{20.2 Technical Power}\label{ch2014martinhairer:technical-power}

The theory of regularity structures is one of the most technically
demanding constructions in modern mathematics. Hairer's ability to
combine algebraic, analytic, and probabilistic ideas into a coherent
whole reflects notable mathematical power.

\subsection{20.3 Generosity}\label{ch2014martinhairer:generosity}

Hairer is known for his generosity in sharing ideas and his commitment
to clear exposition. His lecture notes and monographs are models of
pedagogical clarity, making difficult subjects accessible to students
and researchers.

\begin{center}\rule{0.5\linewidth}{0.5pt}\end{center}

\section{21. Conclusion}\label{ch2014martinhairer:conclusion}

Martin Hairer is one of the influential mathematicians working at
the interface of analysis, probability, and mathematical physics. In a
career spanning two decades, he has: - Solved the KPZ equation, a
central problem in statistical mechanics - Created the theory of
regularity structures, a general framework for singular SPDEs -
Developed the theory of controlled rough paths - Proved ergodicity and
hypoellipticity for infinite-dimensional stochastic systems -
Constructed the dynamical \(\Phi^4_3\) model, a milestone in rigorous
quantum field theory - Trained a generation of researchers in stochastic
analysis

The Fields Medal citation --- \emph{``for outstanding contributions to
the theory of stochastic partial differential equations, and in
particular for the creation of a theory of regularity structures for
such equations''} --- captures the essence of his contributions. The
theory of regularity structures is not merely a solution to a specific
set of problems but a new way of thinking about stochastic dynamics,
with implications that will continue to unfold for decades to come.

\begin{center}\rule{0.5\linewidth}{0.5pt}\end{center}

\section{22. References and Further
Reading}\label{ch2014martinhairer:references-and-further-reading}

\subsection{Official Sources}\label{ch2014martinhairer:official-sources}

\begin{itemize}
\tightlist
\item
 IMU Fields Medal 2014 Announcement:
 https://www.mathunion.org/imu-awards/fields-medal/fields-medals-2014
\item
 Fields Medal Citation for Martin Hairer:
 https://www.mathunion.org/fileadmin/IMU/Prizes/Fields/2014/news\_release\_hairer.pdf
\end{itemize}

\subsection{Biographical and Overview
Articles}\label{ch2014martinhairer:biographical-and-overview-articles}

\begin{itemize}
\tightlist
\item
 Klarreich, E. ``A Mathematician Who Deciphers the Randomness of
 Nature.'' \emph{Quanta Magazine}, August 13, 2014.
\item
 Ball, P. ``The Sound of Mathematics.'' \emph{Nature}, August 2014.
\item
 Royal Society biographical notes.
\end{itemize}

\subsection{Key Technical References}\label{ch2014martinhairer:key-technical-references}

\begin{itemize}
\tightlist
\item
 Hairer, M. ``A theory of regularity structures.'' \emph{Inventiones
 Mathematicae} 198 (2014), 269--504.
\item
 Hairer, M. ``Solving the KPZ equation.'' \emph{Annals of Mathematics}
 178 (2013), 559--664.
\item
 Hairer, M. and Mattingly, J. ``Ergodicity for the stochastic
 Navier-Stokes equations with degenerate noise.'' \emph{Annals of
 Mathematics} 164 (2006), 993--1032.
\item
 Friz, P. and Hairer, M. \emph{A Course on Rough Paths}. Springer, 2014
 (2nd ed.~2020).
\item
 Hairer, M. ``Introduction to Regularity Structures.'' \emph{Brazilian
 Journal of Probability and Statistics} 29 (2015), 175--210.
\end{itemize}
